\documentclass{article}  % Define la clase del documento.

% Paquetes de idioma y codificación
\usepackage[utf8]{inputenc}
\usepackage[T1]{fontenc}
\usepackage[spanish]{babel}  % Ajusta el idioma del documento a español.
\addto\captionsspanish{
  \renewcommand{\figurename}{Figura}
  \renewcommand{\tablename}{Tabla}
}
\usepackage{multirow}

\usepackage{tabularx}  % Permite la creación de tablas con ancho ajustable.

\usepackage{caption}
\usepackage{subcaption}

% Paquete de geometría para configurar márgenes y tamaño de papel
\usepackage[letterpaper, margin=3cm]{geometry}

% Paquetes de tipografía
\usepackage{mathptmx}    % Usa Times New Roman como fuente.
\usepackage{microtype}   % Mejora la justificación del texto.

% Paquetes para manejo de colores y gráficos
\usepackage{xcolor}      % Define y utiliza colores.
\usepackage{graphicx}    % Permite la inserción de imágenes.
\usepackage{tikz}        % Creación de gráficos vectoriales.

% Configuración de enlaces y referencias cruzadas


\usepackage{media9} % Permite la inserción de multimedia.

% Paquetes para la mejora visual de tablas y figuras
\usepackage{booktabs}    % Para tablas de alta calidad.
\usepackage{float}       % Controla la posición de figuras y tablas.

% Paquete para la personalización de códigos fuente
\usepackage{listings}
\lstset{
    literate=
    {á}{{\'a}}1 {é}{{\'e}}1 {í}{{\'i}}1 {ó}{{\'o}}1 {ú}{{\'u}}1
    {Á}{{\'A}}1 {É}{{\'E}}1 {Í}{{\'I}}1 {Ó}{{\'O}}1 {Ú}{{\'U}}1
    {ñ}{{\~n}}1 {Ñ}{{\~N}}1 {ü}{{\"u}}1 {Ü}{{\"U}}1,
    backgroundcolor=\color{backcolour},
    commentstyle=\color{codegreen},
    keywordstyle=\color{codepurple},
    numberstyle=\tiny\color{codegray},
    stringstyle=\color{red},
    basicstyle=\ttfamily\small,
    breakatwhitespace=false,
    breaklines=true,
    captionpos=b,
    keepspaces=true,
    numbers=left,
    numbersep=5pt,
    showspaces=false,
    showstringspaces=false,
    showtabs=false,
    tabsize=2,
    language=TeX,
    morecomment=[l]\#,
    frame=single,
    rulecolor=\color{black}
}

% Definición de colores al estilo Visual Studio Code
\definecolor{darkblue}{rgb}{0.0, 0.0, 0.55}  % Enlaces
\definecolor{codegreen}{rgb}{0.25, 0.49, 0.48}  % Comentarios
\definecolor{codegray}{rgb}{0.5, 0.5, 0.5}  % Números y anotaciones
\definecolor{codepurple}{rgb}{0.58, 0, 0.82}  % Palabras clave
\definecolor{backcolour}{rgb}{0.95, 0.95, 0.92}  % Fondo de código

% Configuraciones de párrafo y matemáticas
\usepackage{amsmath}
\usepackage{parskip}    % Espaciado entre párrafos.
\usepackage{ragged2e}   % Justificación mejorada.
\usepackage{multicol}
\usepackage{multirow}   % Para multirow
\usepackage{float}      % Para la opción [H] en las tablas/figuras


% Configuración de secciones y encabezados
\usepackage{titlesec}
\titleclass{\part}{top} % Make part like a class
\titleformat{\part}[display]
  {\normalfont\huge\bfseries\centering}{\thepart}{40pt}{\Huge}
\titlespacing*{\part}{0pt}{-60pt}{10pt}
\titleformat{\part}
  {\normalfont\huge\bfseries}{}{0pt}{}

% Asegúrate de usar esto para mantener el estilo en las páginas de las partes
\titleformat{\part}[display]
  {\normalfont\huge\bfseries}{}{0pt}{}
  [\thispagestyle{fancy}] % Aplica el estilo fancy a las páginas de las partes

% Configuración de encabezados y pies de página personalizados
\usepackage{fancyhdr}
\pagestyle{fancy}
\fancyhf{}
\fancyhead[L]{\raisebox{0.20cm}{\textbf{Tecnología del Hormigón}}}
\fancyhead[R]{\raisebox{0.1cm}{\includegraphics[width=0.25\linewidth]{LOGO_UNIVERSIDAD.jpg}}}
\fancyhead[C]{\rule{\textwidth}{0.6pt}}
\fancyfoot[C]{\rule{\textwidth}{0.6pt}}
\fancyfoot[R]{\raisebox{-1.5\baselineskip}{\thepage}}
\renewcommand{\headrulewidth}{0pt}
\renewcommand{\footrulewidth}{0pt}

% Configuración avanzada de geometría
\geometry{
  top=3.5cm, % Aumenta el espacio en la parte superior para subir el encabezado
  bottom=2.5cm,
  headheight=2.5cm % Aumenta la altura del encabezado si es necesario
}

% Configuracion de bibliografia
\usepackage{natbib}
\usepackage{hyperref}
\hypersetup{
    colorlinks   = true,
    linkcolor    = darkblue,
    citecolor    = black,
    filecolor    = blue,
    urlcolor     = blue
}

\begin{document}
%----------------------------------------------------------------------------------------
% PORTADA
%----------------------------------------------------------------------------------------
\begin{titlepage}%Inicio de la carátula, solo modificar los datos necesarios
\newcommand{\HRule}{\rule{\linewidth}{0.5mm}} 
\center 
%----------------------------------------------------------------------------------------
%	ENCABEZADO
%----------------------------------------------------------------------------------------
\includegraphics[width=10cm]{LOGO_UNIVERSIDAD.jpg}\\ % Si esta plantilla se copio correctamente, va a llevar la imagen del logo de la facultad.OBS: Es necesario incluir el paquete: graphicx
\vspace{3cm}
%----------------------------------------------------------------------------------------
%	SECCION DEL TITULO
%----------------------------------------------------------------------------------------
\HRule \\[0.4cm]
{ \huge \bfseries Tecnología del Hormigón}\\[0.4cm] % Titulo del documento
{ \huge \bfseries Taller 1}\\[0.4cm] % Titulo del documento
\HRule \\[1.5cm]
 \vspace{5cm}
%----------------------------------------------------------------------------------------
%	SECCION DEL AUTOR
%----------------------------------------------------------------------------------------
\begin{flushright}
  { \textbf{Profesor:}\\
  Alvaro Paul\\
  \textbf{Ayudante:}\\
  Felipe Ronda\\
  \textbf{Alumnos:} \\
  Felipe Vicencio\\
  Lukas Wolff\\
}
\end{flushright}
\vspace{1cm}
%----------------------------------------------------------------------------------------
%	SECCION DE LA FECHA
%----------------------------------------------------------------------------------------
{\large \textbf{\today}}\\[2cm] % El comando \today coloca la fecha del dia, y esto se actualiza con cada compilacion, en caso de querer tener una fecha estatica, reemplazar el \today por la fecha deseada
\end{titlepage}

\newpage
\thispagestyle{empty} % Deshabilita el número de página en la página del índice

%----------------------------------------------------------------------------------------
%  INDICE
%----------------------------------------------------------------------------------------
%\newpage
%\thispagestyle{empty} % Deshabilita el número de página en la página del índice
%\tableofcontents
%\thispagestyle{plain} % Deshabilita el encabezado en la página del índice
%\thispagestyle{empty} % Deshabilita el número de página en la página del índice
%\newpage

%\newpage
%\thispagestyle{empty}
%\listoffigures 
%\thispagestyle{plain} % Deshabilita el encabezado en la página del índice %
%\thispagestyle{empty}
%----------------------------------------------------------------------------------------
%ACÁ EMPIEZA EL INFORME
\setcounter{page}{1}
%----------------------------------------------------------------------------------------

\section{Introducción}



\section{Desarrollo}

\subsection{Metodología}

\subsubsection{Análisis estadístico de la resistencia mecánica}

En primer lugar, se realizó un análisis estadístico de la resistencia mecánica, basado en la norma \cite{NCh1998}, cuyo objetivo es determinar la conformidad de los resultados de la resistencia real a compresión respecto de la especificada.

Para poder aplicar la norma, se deben cumplir los siguientes requisitos básicos:

\begin{itemize}
    \item Toma de muestras al azar.
    \item Toma de muestras entre el 10\% y el 90\% de la descarga del camión mixer.
    \item Tamaño de muestra superior a 1,5 veces el necesario.
    \item Generar al menos 3 probetas por muestra: una ensayada a 7 días y las restantes a 28 días.
\end{itemize}

La Tabla \ref{tab:muestreo} detalla el número y volumen de muestras según la procedencia del hormigón.

\begin{table}[H]
\centering
\begin{tabular}{|l|c|c|}
\hline
\textbf{Procedencia del hormigón} & \multicolumn{2}{c|}{\textbf{Volumen de hormigón de la obra, m$^3$}} \\ \hline
                                  & $> 250$ & $\leq 250$ \\ \hline
\textbf{Fabricado en obra}        &         &            \\ \hline
Volumen máximo de hormigón por muestra & 100     & 50        \\ \hline
Número mínimo de muestras              & 5       & 3         \\ \hline
\textbf{De central hormigonera}   &         &            \\ \hline
Volumen máximo de hormigón por muestra & 150     & 75        \\ \hline
Número mínimo de muestras              & 5       & 3         \\ \hline
\end{tabular}
\caption{Requisitos de muestreo según la procedencia y el volumen de hormigón.}
\label{tab:muestreo}
\end{table}

Además, el plan de muestreo debe estar especificado en la norma de diseño o en la especificación de la obra. En caso contrario, se debe considerar lo siguiente:

\begin{itemize}
    \item Pavimentos: un testigo cada 1000 m$^2$, con un mínimo de 3.
    \item General: testigos por zona.
\end{itemize}

Si se dispone de un número $N$ mayor a 10 muestras, se puede considerar que la resistencia de cada lote es satisfactoria si se cumple lo siguiente:

\begin{itemize}
    \item Criterio global: $f_3 \geq f_c + k_1$
    \item Criterio individual: $f_i \geq f_0 = f_c - k_2$
\end{itemize}

Donde $f_3$ corresponde a la resistencia media de tres muestras consecutivas a 28 días, $f_c$ es la resistencia especificada, $f_i$ es la resistencia individual de cada muestra y $f_0$ es la resistencia mínima aceptable para cada $f_i$ a 28 días.

Los valores de $k_1$ y $k_2$ se indican en la Tabla \ref{tab:valoresk}.

\begin{table}[H]
\centering
\begin{tabular}{|c|c|c|c|c|c|}
\hline
\textbf{Fracción defectuosa aceptada, \%} &  & \multicolumn{4}{c|}{\textbf{Grado de hormigón}} \\ \hline
 &  & H5 & H10 & H15 & H20 o superior \\ \hline
\multirow{2}{*}{5}  & $k_{1}$ & 0,3 & 0,5 & 0,8 & 1,0 \\ \cline{2-6}
                    & $k_{2}$ & 0,6 & 1,2 & 1,9 & 2,5 \\ \hline
\multirow{2}{*}{10} & $k_{1}$ & 0   & 0   & 0   & 0   \\ \cline{2-6}
                    & $k_{2}$ & 0,9 & 1,7 & 2,6 & 3,5 \\ \hline
\multirow{2}{*}{20} & $k_{1}$ & 0,4 & 0,7 & 1,1 & 1,5 \\ \cline{2-6}
                    & $k_{2}$ & 1,4 & 2,7 & 4,1 & 5,5 \\ \hline
\end{tabular}
\caption{Valores de $k_{1}$ y $k_{2}$ según la fracción defectuosa aceptada y el grado de hormigón.}
\label{tab:valoresk}
\end{table}

La Tabla \ref{tab:conclusiones} resume las conclusiones posibles en función de los antecedentes.

\begin{table}[H]
\centering
\begin{tabular}{|c|c|p{6cm}|p{6cm}|}
\hline
\textbf{Antecedentes} &  & \textbf{Conclusiones} & \textbf{Recomendaciones} \\ \hline
$f_{3} \geq f_{c} + k_{1}$ & $f_{i} \geq f_{0}$ & 
El hormigón cumple la resistencia especificada & \\ \hline

$f_{3} < f_{c} + k_{1}$ & $f_{i} \geq f_{0}$ & 
El hormigón no cumple la resistencia especificada & 
Informar a los proyectistas estructurales y considerar las penalizaciones establecidas en el contrato y sus documentos anexos. \\ \hline

 & $f_{i} < f_{0}$ & 
El hormigón no cumple la resistencia especificada y cada resultado defectuoso debe ser considerado como un riesgo potencial & 
Adoptar las medidas indicadas en A.4 \\ \hline
\end{tabular}
\caption{Conclusiones y recomendaciones en función de los antecedentes de resistencia del hormigón.}
\label{tab:conclusiones}
\end{table}

Adicionalmente, se puede realizar un análisis considerando el total de muestras. El ensayo se considera satisfactorio si se cumple que:

\begin{itemize}
    \item Criterio global: $f_{m} \geq f_{c} + s \cdot t$
    \item Criterio individual: $f_{i} \geq f_{0} = f_c - k_2$
\end{itemize}

Donde $s$ corresponde a la desviación estándar de la muestra, y el valor $t$ puede calcularse a partir de los datos de la Tabla 4 de la norma \cite{NCh1998}.

La Tabla \ref{tab:conclusiones2} muestra las conclusiones obtenidas en este análisis.

\begin{table}[H]
\centering
\begin{tabular}{|c|c|p{6cm}|p{6cm}|}
\hline
\textbf{Antecedentes} &  & \textbf{Conclusiones} & \textbf{Recomendaciones} \\ \hline
$f_{m} \geq f_{c} + s \cdot t$ & $f_{i} \geq f_{0}$ & 
El hormigón cumple la resistencia especificada & \\ \hline

$f_{m} < f_{c} + s \cdot t$ & $f_{i} \geq f_{0}$ & 
El hormigón no cumple la resistencia especificada & 
Informar a los proyectistas estructurales y considerar las penalizaciones establecidas en el contrato y sus documentos anexos. \\ \hline

 & $f_{i} < f_{0}$ & 
El hormigón no cumple la resistencia especificada y cada resultado defectuoso debe ser considerado como un riesgo potencial & 
Adoptar las medidas indicadas en A.4 \\ \hline
\end{tabular}
\caption{Conclusiones y recomendaciones considerando $f_{m}$, $f_{c}$, $s$, $t$ y $f_{0}$.}
\label{tab:conclusiones2}
\end{table}

Finalmente, si se dispone de por lo menos 10 muestras distitnas, con dos o mas probetas por cada una, se puede realizar una evaluacion del nivel de control de los ensayos.

En primer lugar es nesesario calcular el intervalo promedio R:

\begin{equation}
    \bar{R} = \frac{\sum R_i}{N}
\end{equation}

Donde:

\begin{itemize}
    \item $R_i = f_{max,i} - f_{min,i}$ es el rango de cada muestra.
    \item $N$ es el número total de muestras.
\end{itemize}

A continuación, se debe calcular la desviacion normal de los ensayos:

\begin{equation}
    s_1 = \frac{\bar{R}}{d_2}
\end{equation}

Para el valor de d2, considerar los datos de la tabla \ref{tab:d2}

\begin{table}[H]
\centering
\begin{tabular}{|c|c|}
\hline
\textbf{Número de probetas, $n_{0}$} & \textbf{$d_{2}$} \\ \hline
2 & 1,128 \\ \hline
3 & 1,693 \\ \hline
4 & 2,059 \\ \hline
5 & 2,326 \\ \hline
6 & 2,534 \\ \hline
\end{tabular}
\caption{Valores de $d_{2}$ en función del número de probetas $n_{0}$.}
\label{tab:d2}
\end{table}

Finalmente se calcula el coeficiente de vaiacion de ensayo:

\begin{equation}
    V_1 = \frac{s_1}{f_m} \cdot 100\%
\end{equation}

Donde la tabla \ref{tab:conclusiones3} muestra las conclusiones y recomendaciones en función del coeficiente de variación obtenido.

\begin{table}[H]
\centering
\begin{tabular}{|c|c|}
\hline
\textbf{$V_{1}$ (\%)} & \textbf{Nivel de control de ensayos} \\ \hline
$0 \leq V_{1} \leq 3,0$ & Excelente \\ \hline
$3,0 < V_{1} \leq 4,0$ & Muy bueno \\ \hline
$4,0 < V_{1} \leq 5,0$ & Bueno \\ \hline
$5,0 < V_{1} \leq 6,0$ & Aceptable \\ \hline
$6,0 < V_{1}$ & Deficiente \\ \hline
\end{tabular}
\caption{Clasificación del nivel de control de ensayos según el valor de $V_{1}$.}
\label{tab:conclusiones3}
\end{table}


\subsubsection{Analisis de Testigos}

La norma \cite{NCh1171-1-2001} establece los requicitos de analisis para la extraccion de testigos. Se deben cumplir los siguientes requerimientos:

\begin{itemize}
    \item El diametro debe ser mayor que el doble del diametro de maximo agregado
    \item Altura tal que la esbeltez sea entre 1.0 y 2.0
    \item La resistencia de extraccion debe ser mayor a 8 MPa y mayor a 14 dias.
    \item Ademas, los limites fisicos de extraccion son especificados por la norma.
\end{itemize}












\subsection{Resultados}

Los datos de entrada son un hormigon G55 con un porcentaje de defectuoso del 20\%, ademas, se tienen las siguientes muestras y probetas (tabla \ref{tab:datosentrada}).

\begin{table}[H]
\centering
\begin{tabular}{|c|c|c|c|}
\hline
\textbf{Correlativo} & \textbf{7 días [MPa]} & \textbf{28 días [MPa]} & \textbf{28 días [MPa]} \\ \hline
1  & 41.45 & 52.8  & 54.6  \\ \hline
2  & 40.2  & 57.11 & 56.6  \\ \hline
3  & 43.5  & 52.35 & 55.18 \\ \hline
4  & 41    & 55.25 & 54.2  \\ \hline
5  & 40    & 54.48 & 53.82 \\ \hline
6  & 41    & 55.3  & 52.08 \\ \hline
7  & 43.5  & 51    & 53.82 \\ \hline
8  & 41    & 48.4  & 49.9  \\ \hline
9  & 40    & 49.5  & 45.6  \\ \hline
10 & 40.2  & 52.56 & 47.88 \\ \hline
11 & 43.5  & 57.11 & 56.6  \\ \hline
12 & 40    & 52.35 & 52.08 \\ \hline
13 & 41    & 55.25 & 53.82 \\ \hline
14 & 43.5  & 54.48 & 53.82 \\ \hline
15 & 41    & 55.3  & 57.79 \\ \hline
16 & 40    & 58.15 & 56.6  \\ \hline
17 & 40.2  & 55.25 & 55.18 \\ \hline
18 & 43.5  & 52.56 & 53.82 \\ \hline
19 & 41    & 48.4  & 45.68 \\ \hline
20 & 40    & 49.6  & 47.7  \\ \hline
\end{tabular}
\caption{Resultados de resistencia a compresión de probetas a 7 y 28 días.}
\label{tab:datosentrada}
\end{table}

Luego, se obtienen los siguientes resultados: 

\begin{table}[H]
\centering
\begin{tabular}{|c|c|c|}
\hline
\textbf{Probeta} & \textbf{Cumple $f_c - k_2$} & \textbf{Cumple $f_3 \geq f_c + k_1$} \\ \hline
1  & Sí & Sí \\ \hline
2  & Sí & Sí \\ \hline
3  & Sí & Sí \\ \hline
4  & Sí & Sí \\ \hline
5  & Sí & No \\ \hline
6  & Sí & No \\ \hline
7  & Sí & No \\ \hline
8  & No & No \\ \hline
9  & No & No \\ \hline
10 & Sí & No \\ \hline
11 & Sí & Sí \\ \hline
12 & Sí & Sí \\ \hline
13 & Sí & Sí \\ \hline
14 & Sí & Sí \\ \hline
15 & Sí & Sí \\ \hline
16 & Sí & Sí \\ \hline
17 & Sí & No \\ \hline
18 & Sí & No \\ \hline
19 & No & - \\ \hline
20 & No & - \\ \hline
\end{tabular}
\caption{Resultados de cumplimiento de las condiciones estadísticas de resistencia del hormigón.}
\label{tab:cumplimiento}
\end{table}

Ademas, el criterio $f_m \geq f_c + st$ no se cumple. De esta forma, se puede concluir que es nesesario adoptar las medidas indicadas en A.4. de la norma \cite{NCh1998}.

Se concluye que:

\begin{itemize}
    \item Existe riesgo estructural, y el hormigon afectado debe ser sometido a analisis
    \item Se debe comprobar la validex del ensayo, ademas de extraer testigos de hormigon endurecido, identificar zonas comprometidas y demas.
    \item \textbf{CONTINUAAAAAR}
\end{itemize}

Adicionalmente, se calcula el nivel de control del ensayo, donde se concluye que su calidad es excelente.

Al modificar el porcentaje de defectuoso, no se genera ningun cambio al aumentarlo, ya que la dosificacion entregada se encuentra en el limite superior (20\%), de manera contraria, si se reduce tal porcentaje, se reduce el numero de muetras aprobadas, ya que los factores k aumentan, incrementando la diferencia de la desigualdad. Por ejemplo, incialmente se tienen 26 muestraqs aprobadas considerando los dos criterior individuales, pero si el porcentaje de defectuso cambia a 5\%, se tienen 15 probetas aprobadas.

\section{Conclusiones}

En conclusión, este taller buscó realizar un estudio del análisis estadístico de la resistencia mecánica del hormigón extraído de obras, aplicando las normas NCh1998 y NCh1171-1-2001 para evaluar la conformidad de los resultados obtenidos. A partir de los datos de entrada, se determinó que el hormigón no cumple la resistencia especificada, lo que implica un riesgo estructural y la necesidad de adoptar medidas correctivas. Asimismo, se evaluó el nivel de control de los ensayos, obteniéndose un nivel excelente, lo que indica un adecuado proceso de muestreo y una baja variabilidad entre probetas.

Además, se discutió qué normas prescriptivas y de desempeño pueden aplicarse en distintas situaciones, en conexión con el desarrollo del taller anterior, en el que se llevó a cabo una optimización granulométrica para el hormigón solicitado.

Finalmente, se generaron diversas discusiones sobre los efectos de distintos factores en la mezcla de hormigón, como la adición de productos químicos y la porosidad, entre otros, lo que permitió no solo analizar las muestras por sus resultados mecánicos, sino también alcanzar una comprensión más amplia de los múltiples elementos que pueden estar involucrados.

\include{anexo}

\bibliographystyle{unsrtnat}  % o plainnat
\bibliography{ref}  % Nombre del archivo .bib

\end{document}